\section{Introduction}
Type theory is a logical formalism used extensively
in the study and design of programming languages to define the
semantics and behavior of deductive systems.
According to~\cite{harpercourse},
``The central organizing principle of language design is the
identification of language features with types.
The theory of programming languages, therefore, reduces to the
theory of types.''
Type theory ranges over several formal systems,
but for this paper, type theory refers to
the design, analysis and study of \emph{type systems}.
Due to the wide variety of usage of the phrase
``type system'', it lacks a standard definition.
According to~\cite{typesystems-wikipedia}:
``A type system is a collection of rules that assign a property
called a type to the various language constructs---such as variables, expressions, functions or modules---that a computer program is composed of.''
A type system formally defines many aspects of a programming language
such as in the following far-from-exhaustive list:
\begin{enumerate}
\item Which expressions are allowed in the language;
specifically, what a well-formed or well-\emph{typed} program in the
language is.
Vaguely, well-typed programs are programs that entail certain
properties such as the absence of certain program errors~\cite{Pierce:2002:TPL:509043}.
\item How program abstractions and components of
large systems can be tied together.
\item How expressions in the language are evaluated.
What is the order of evaluation?
Are expressions evaluated eagerly
(every expression in a program is computed immediately)
or lazily
(expressions are not evaluated until the program
requires their values)?
\end{enumerate}
Type systems model complex language features
and enable one to prove properties about languages.
Type systems are rigorous enough to be
used as a language specification for a compiler
writer to implement the language;
the implementation of a programming language
can follow from its type system.

A type system is specified by a set of inference rules
that define a programming language.
These inference rules are partitioned into two categories.
Rules defining the types of the \emph{terms} or expressions in
the language are the \emph{static semantics}.
Static semantics \emph{inductively} define
a relation between expressions and types.
\emph{Dynamic semantics} or operational semantics inductively
define how to evaluate expressions in the language.
Specifically, dynamic semantics define
a transition system between expressions, where
the \emph{values} that expressions evaluate to are
the \emph{final states} of the system.

Type systems are best explained with an example.
The following sections define a programming language
that we call $\minilang$.
The \emph{grammar} of $\minilang$ is defined in
Section~\ref{sec:grammar}.
Its static and dynamic semantics are
given in Sections~\ref{sec:static}
and~\ref{sec:dynamic}, respectively.
We prove an important property, \emph{type preservation},
for $\minilang$ in Section~\ref{sec:preservation}.
We show how to prove another important property,
\emph{progress}, by proving it
for one type of expression in $\minilang$
in Section~\ref{sec:progress}.
The proof of progress for all other expressions
is similar, so we skip those cases for brevity.

Proofs of language properties are important
but long, tedious, and error-prone
because there are many cases to reason about.
As a result, \emph{proof assistants}
have been developed for proving language theory.
These software tools provide a language
for writing properties and proofs of properties.
These tools can find mistakes in proofs
and verify proof correctness.
We cover one proof assistant, Twelf~\cite{Twelf},
in Section~\ref{sec:twelf}.
All of the Twelf code presented in this paper can
be found in~\cite{twelf-tutorial}.
Slides accompany the paper: one deck presents
$\minilang$ and its type safety proofs~\cite{typetheory-slides},
and another presents its Twelf encoding~\cite{twelf-slides}.
Section~\ref{sec:summary} concludes with a summary.
